% BHU PYQ -- Manually transcribed & error-corrected
% Paper: GLB-201: Elements of Mineralogy and Crystallography
% Exam: B.Sc. (Hons.) Semester II Examination, 2013-14

\documentclass[11pt,a4paper]{article}
\usepackage[a4paper,top=2.5cm,bottom=2.5cm,left=2.8cm,right=2.8cm,headheight=14pt]{geometry}
\usepackage{amsmath,amssymb}
\usepackage[T1]{fontenc}
\usepackage[utf8]{inputenc}
\usepackage{lmodern,microtype,fancyhdr,enumitem,hyperref,lastpage,parskip}

\hypersetup{colorlinks=true,urlcolor=black,linkcolor=black,
  pdftitle={GLB-201 Elements of Mineralogy and Crystallography -- BHU B.Sc. Sem II 2013-14}}
\pagestyle{fancy}\fancyhf{}
\renewcommand{\headrulewidth}{0.4pt}\renewcommand{\footrulewidth}{0.4pt}
\fancyhead[L]{\small   \textit{Banaras Hindu University}}
\fancyhead[R]{\small   \textit{GLB-201: Elements of Mineralogy \& Crystallography}}
\fancyfoot[C]{\small Page  \thepage\ of \pageref{LastPage}}


\newlist{parts}{enumerate}{2}
\setlist[parts,1]{label=(\alph*),leftmargin=*,itemsep=5pt,topsep=3pt}

\newlist{subparts}{enumerate}{2}
\setlist[subparts,1]{label=(\roman*),leftmargin=*,itemsep=3pt,topsep=2pt}

\newcommand{\pts}[1]{\hfill   \textit{[#1]}}


\begin{document}

\begin{center}
  {\large\bfseries BANARAS HINDU UNIVERSITY}\\[2pt]
  {
\normalsize B.Sc. (Hons.) --- Semester II Examination, 2013-14}\\[6pt]
  \rule{0.8\linewidth}{0.5pt}\\[6pt]
  {\large\bfseries Paper: GLB-201 --- Elements of Mineralogy and Crystallography}\\[4pt]
  {
\normalsize Time: 3 Hours \hfill Maximum Marks: 70}\\[2pt]
  {\small Ref. No.: BS/Sem II/16}
\end{center}
\rule{\linewidth}{0.4pt}
\smallskip



\noindent   \textit{Instructions: Answer   \textbf{five} questions in all, selecting at least   \textbf{two} each from Sections A and B, and Section-C which is compulsory.}
\bigskip

\bigskip


\noindent {\large\bfseries SECTION --- A}
\rule{\linewidth}{0.4pt}\smallskip




\noindent   \textbf{Question 1.} What is magma? Explain the magmatic process of mineral formation. \pts{14}

\medskip


\noindent   \textbf{Question 2.} Explain with neat diagrams, the determination of Miller indices in different crystal systems. \pts{14}

\medskip


\noindent   \textbf{Question 3.} With a neat diagram, describe different parts and their functions in a petrological microscope. \pts{14}

\bigskip


\noindent {\large\bfseries SECTION --- B}
\rule{\linewidth}{0.4pt}\smallskip




\noindent   \textbf{Question 4.} Answer the following parts: \pts{14}

\begin{parts}
  \item Which among the following is a mineral? \pts{7}
  
\begin{subparts}
    \item Wood
    \item Quartz
  \end{subparts}
  Give proper justification in support of your answer.
  
  \item What do you understand by the hardness of a mineral? Explain Mohs' scale of hardness. \pts{7}
\end{parts}

\medskip


\noindent   \textbf{Question 5.} Write the chemical composition and physical properties of   \textbf{any four} of the following minerals underlining the diagnostic ones: \pts{$4  \times 3\frac{1}{2} = 14$}

\begin{parts}
  \item Hematite
  \item Gypsum
  \item Calcite
  \item Kyanite
  \item Bauxite
\end{parts}

\medskip


\noindent   \textbf{Question 6.} Write short notes on   \textbf{any four} of the following: \pts{$4  \times 3\frac{1}{2} = 14$}

\begin{parts}
  \item Centre of symmetry
  \item First Law of Crystallography
  \item Birefringence
  \item Pleochroism
  \item Twinkling
\end{parts}

\bigskip


\noindent {\large\bfseries SECTION --- C}
\rule{\linewidth}{0.4pt}\smallskip




\noindent   \textbf{Question 7.} Answer the following questions:

\begin{parts}
  \item Whether the following statement is True or False --- ``On a streak plate the streak of topaz is yellow in color.'' Justify your answer with reasons. \pts{2}
  \item Arrange the following minerals in decreasing order of resistance to weathering: \pts{2}
  
\begin{center}
    Hornblende; K-feldspar; Olivine; Quartz.
  \end{center}
  \item Coal and petroleum are products of \underline{\hspace{4cm}} process. (Fill in the blank) \pts{1}
  \item Differentiate between the terms `crystal' and `crystalline'. \pts{2}
  \item Name the crystallographic axes of calcite. \pts{1}
  \item What is a prism face in a crystal? \pts{1}
  \item Write the crystal systems of calcite and albite. \pts{2}
  \item Draw the ray velocity surface for an anisotropic mineral. \pts{1}
  \item What is a parametral face of a crystal? \pts{1}
  \item Give one example of an isotropic mineral. \pts{1}
\end{parts}

\vfill\rule{\linewidth}{0.4pt}

\begin{center}\small
\href{https://bhu-exam.vercel.app/}{Click here to visit our website}\\[4pt]
\href{https://me-aryan.vercel.app/}{Click here to visit our contributor}\\[4pt]
\href{https://quashan-me.vercel.app/}{Click here to visit our contributor}
\end{center}
\end{document}
